% ======================================================================
%  Mémoire de fin d'année — PGE ING3
%  Conception d'une plateforme de formation autonome alimentée par l'IA
%  Auteur : ISSAD Nassim — Entreprise : Audit Énergie
%
%  FICHIER UNIQUE ET AUTONOME (bibliographie incluse).
%  Compilation : pdflatex memoire-complet  (x2 pour le sommaire)
% ======================================================================
\documentclass[12pt,a4paper]{article}

\usepackage[utf8]{inputenc}
\usepackage[T1]{fontenc}
\usepackage{lmodern}
\usepackage[french]{babel}

\usepackage[a4paper,margin=2.5cm]{geometry}
\usepackage{setspace}
\onehalfspacing
\setlength{\parskip}{0.6em}
\setlength{\parindent}{0pt}

\usepackage{array}
\usepackage{graphicx}
\usepackage{float}

\usepackage[dvipsnames]{xcolor}
\usepackage[hidelinks]{hyperref}
\hypersetup{
  pdftitle={Conception d'une plateforme de formation autonome alimentée par l'IA},
  pdfauthor={ISSAD Nassim}
}

\usepackage{fancyhdr}
\pagestyle{fancy}
\fancyhf{}
\fancyhead[L]{\small\itshape Plateforme de formation autonome}
\fancyhead[R]{\small Audit Énergie}
\fancyfoot[C]{\thepage}
\renewcommand{\headrulewidth}{0.3pt}

\begin{document}

% ====================== PAGE DE GARDE ===============================
\begin{titlepage}
\thispagestyle{empty}
\centering

{\scshape\large EFREI Paris \par}
\vspace{0.3cm}
{\scshape Mémoire de fin d'année --- PGE ING3 --- 2025-2026 \par}
\vspace{0.2cm}
\rule{\linewidth}{0.5pt}

\vfill

{\Huge\bfseries Conception d'une plateforme de\\[0.3cm]
formation autonome alimentée par\\[0.3cm]
l'intelligence artificielle \par}

\vspace{0.8cm}
{\large\itshape Pipeline de génération de contenus pédagogiques et\\
assistant conversationnel contextualisé (RAG)\par}

\vfill

\begin{tabular}{@{}ll@{}}
\textbf{Auteur} & ISSAD Nassim \\[0.3cm]
\textbf{Promotion} & ING3 --- M2-BDML \\[0.3cm]
\textbf{Tuteur entreprise} & M. Mehdi ROUSSELLE \\[0.3cm]
\textbf{Tuteur pédagogique} & M. Moulay AHANSAL \\[0.3cm]
\textbf{Chargé de Mission Career Center} & \textit{à compléter} \\[0.3cm]
\textbf{Entreprise} & Audit Énergie \\
\end{tabular}

\vfill

\rule{\linewidth}{0.5pt}
{\small Diffusion soumise à l'accord de l'entreprise --- \textit{mention de confidentialité à compléter}\par}
\vspace{0.3cm}
{\small Année académique 2025--2026\par}

\end{titlepage}

% ====================== RÉSUMÉ ======================================
\clearpage
\pagenumbering{roman}
\phantomsection
\addcontentsline{toc}{section}{Résumé}
\section*{Résumé}

Ce mémoire présente la conception d'une plateforme de formation autonome alimentée par
l'intelligence artificielle, développée au sein de l'entreprise Audit Énergie. Initialement
positionnée sur le courtage en énergie, l'entreprise s'est récemment structurée comme un
organisme de formation externe préparant des apprenants à des titres professionnels. Pour
une petite structure, ce développement se heurte à une contrainte forte : le coût d'un
formateur en direct rend le modèle difficilement viable, tandis que les cours simplement
préenregistrés appauvrissent l'accompagnement pédagogique.

Le travail propose de répondre à ce problème par une plateforme capable d'automatiser la
production, la diffusion et l'accompagnement des cours. Deux contributions principales sont
étudiées. La première est une \emph{pipeline} de génération de contenus pédagogiques, qui
transforme un référentiel officiel en cours structurés, oralisables et calibrés en durée. Plutôt
qu'une génération directe par un prompt unique, cette pipeline s'appuie sur un plan
intermédiaire explicite, une génération section par section, des contrôles qualité ciblés et la
conservation d'artefacts intermédiaires garantissant la traçabilité. La seconde contribution est
un assistant pédagogique de type RAG (\emph{Retrieval-Augmented Generation}), fondé sur
Azure AI Search et Azure OpenAI, qui répond aux questions des élèves à partir du contenu réel
de la formation plutôt que d'une connaissance générale non contrôlée.

Le mémoire situe ces choix dans l'état de l'art (prompt engineering, fine-tuning, RAG,
génération structurée, reviews automatiques) et discute leurs limites. Il insiste sur une
exigence centrale : encadrer l'IA générative pour la rendre fiable, auditable et mesurable dans
un usage réel de production. Sont enfin présentées les conditions d'application de l'approche,
ses limites et les indicateurs d'évaluation associés, qu'il s'agisse de l'efficacité opérationnelle,
de la qualité pédagogique ou de la pertinence des réponses de l'assistant.

\vspace{0.4cm}
\noindent\textbf{Mots-clés :} intelligence artificielle générative, formation à distance,
pipeline de génération de contenus, synthèse vocale, RAG, Azure AI Search, traçabilité,
automatisation pédagogique.

% ====================== SOMMAIRE ====================================
\clearpage
\phantomsection
\addcontentsline{toc}{section}{Sommaire}
\tableofcontents

% ====================== CORPS DU MÉMOIRE ============================
\clearpage
\pagenumbering{arabic}

\phantomsection
\addcontentsline{toc}{section}{Introduction générale}
\section*{Introduction générale}

\subsection*{Contexte}

Audit Énergie est une entreprise initialement positionnée sur le courtage en énergie, dont
le métier consiste à accompagner des professionnels dans la renégociation de leurs contrats.
Cette activité lui a donné une connaissance fine des métiers de la vente à distance. C'est
cette expertise qui l'a conduite, plus récemment, vers la formation professionnelle :
d'abord pour ses propres téléprospecteurs, puis, depuis environ un an, sous la forme d'un
organisme de formation externe préparant des apprenants à des titres professionnels reconnus
par l'État, comme le titre de Conseiller Relation Client à Distance ou celui d'Employé
Commercial. C'est cette activité de formation qui constitue le cœur de ce mémoire.

Ce changement d'échelle a fait apparaître une contrainte forte. Pour une petite structure,
faire appel à des formateurs en direct pour assurer les cours en ligne pèse lourdement sur
des marges déjà limitées. Une première réponse a consisté à diffuser des cours audio préparés
à l'avance plutôt qu'animés en direct ; mais ce modèle restait dépendant d'un travail humain
important, et un cours préenregistré ne permet pas de répondre aux questions des élèves.
L'idée s'est donc imposée de concevoir une plateforme de formation réellement autonome,
capable à la fois de produire et de diffuser les contenus, et d'accompagner les apprenants.

\subsection*{Problématique, objectifs et contributions}

Le projet étudié dans ce mémoire est précisément cette plateforme de formation autonome
alimentée par l'intelligence artificielle. Elle ne doit pas seulement diffuser des fichiers
audio, mais générer les contenus, les transformer en audio, produire des supports et des
exercices, et accompagner les élèves grâce à un assistant s'appuyant sur les documents réels
de la formation. La problématique peut alors être formulée ainsi :

\begin{quote}
\emph{Comment concevoir une plateforme de formation autonome alimentée par
l'intelligence artificielle, capable de produire, diffuser et accompagner des contenus
pédagogiques, tout en garantissant leur qualité, leur traçabilité et leur pertinence pour
les apprenants ?}
\end{quote}

Cette problématique se décline en quatre objectifs : automatiser la production des cours, en
passant d'un référentiel officiel à des journées pédagogiques complètes ; transformer ces
contenus en audios exploitables, à un coût inférieur à une production humaine répétée ;
enrichir la formation par des supports visuels et des exercices générés automatiquement ;
et permettre aux apprenants d'interroger un assistant capable de répondre à partir du contenu
réel du cours, et non comme un chatbot généraliste.

Pour y répondre, le travail apporte deux contributions principales. La première est une
\emph{pipeline} de génération de contenus pédagogiques, qui a évolué d'une génération
directe de texte vers une chaîne structurée, fondée sur un plan pédagogique explicite, une
génération section par section, des contrôles qualité ciblés et la conservation d'artefacts
intermédiaires garantissant la traçabilité. La seconde est un assistant pédagogique de type
RAG (\emph{Retrieval-Augmented Generation}), qui ancre ses réponses dans les documents
indexés de la formation pour dépasser les limites d'un modèle de langage utilisé seul.

La méthode suivie est itérative : chaque solution essayée (intervention humaine,
enregistrement différé, synthèse vocale, génération directe par modèle de langage, puis
pipeline structurée) a révélé une nouvelle contrainte — coût, qualité audio, longueur des
textes, cohérence pédagogique, conformité, accompagnement, évaluation. L'enjeu n'est pas
seulement d'ajouter de l'intelligence artificielle à une plateforme existante, mais de
l'encadrer pour la rendre fiable, contrôlable et mesurable. Les résultats sont donc discutés à
l'aide d'indicateurs concrets : réduction des interventions humaines, coût et temps de
production des contenus, et, pour l'assistant, pertinence des passages retrouvés et fidélité des
réponses produites.

\subsection*{Structure du mémoire}

Ce mémoire est structuré en trois parties. La première, \emph{Position du problème}, présente
le contexte de l'entreprise, les limites du fonctionnement initial et la problématique de la
formation autonome. La deuxième, \emph{État de l'art}, étudie les approches existantes :
plateformes de formation en ligne, synthèse vocale, génération de contenus par IA et systèmes
de questions-réponses augmentés par récupération d'information. La troisième, \emph{Étude
originale}, présente la solution développée, ses choix d'architecture, les difficultés
rencontrées, les arbitrages réalisés et les méthodes retenues pour en évaluer les apports.

% ====================== PARTIE 1 — POSITION DU PROBLÈME ======================
\clearpage
\section{Position du problème}

\subsection{D'un courtier en énergie à un organisme de formation}

Audit Énergie s'est d'abord construite autour du courtage en énergie. Son métier consiste
à accompagner des clients professionnels dans la compréhension, la comparaison et la
renégociation de leurs contrats, sur un marché devenu progressivement plus complexe. Au
fil des années, cette activité a forgé une expertise concrète des métiers de la vente à
distance : la prospection, la relation client par téléphone, l'analyse d'un besoin
commercial, l'argumentation et le suivi d'un dossier dans la durée.

Ce positionnement initial éclaire le sujet du mémoire, mais il n'en constitue pas le cœur.
Le projet étudié s'inscrit dans une évolution plus récente : la transformation d'Audit
Énergie en organisme de formation. Cette évolution s'est faite en deux temps. D'abord en
interne, pour former ses propres téléprospecteurs aux fondamentaux de la relation client
et aux spécificités du secteur. Puis, depuis environ un an, de manière beaucoup plus
structurée, avec le développement d'une offre de formation externe. L'entreprise prépare
désormais des apprenants à des titres professionnels reconnus par l'État, notamment le
titre de Conseiller Relation Client à Distance ou celui d'Employé Commercial.

Ce changement d'échelle modifie en profondeur les contraintes de l'entreprise. Former
quelques collaborateurs en interne et faire fonctionner un organisme de formation externe
ne relèvent pas de la même organisation. Dans le premier cas, la transmission peut
s'appuyer sur l'expérience des salariés, des échanges directs et des supports peu
formalisés. Dans le second, il faut produire des parcours reproductibles, préparer des
contenus, assurer une continuité pédagogique, suivre la présence des élèves, proposer des
exercices, répondre aux questions et garantir une qualité suffisante pour que la formation
soit réellement exploitable.

Le sujet du mémoire naît précisément de cette transition. Il ne s'agit pas seulement
d'ajouter une fonctionnalité technique à une plateforme, mais de répondre à une question
d'organisation plus large : comment une petite structure peut-elle développer une activité
de formation à distance sans disposer des moyens humains et financiers d'un organisme de
formation déjà établi ?

\subsection{La formation à distance : une opportunité freinée par son coût}

La formation à distance représente, sur le papier, une opportunité idéale pour une petite
entreprise. Elle permet de toucher des apprenants sans salle physique, de mutualiser des
contenus, de planifier des sessions souples et de réutiliser des supports d'une promotion à
l'autre. Elle semble donc parfaitement compatible avec les contraintes d'une structure
légère.

Cette impression est pourtant incomplète. Une formation à distance n'est rentable que si
ses coûts de production et d'animation restent maîtrisés. Or elle ne se résume pas à
déposer des documents en ligne : il faut concevoir un déroulé pédagogique, produire les
cours, les rendre accessibles, accompagner les apprenants, suivre leur participation et
vérifier leur compréhension. Ces tâches sont répétitives, mais indispensables.

Le modèle le plus naturel aurait été de faire appel à des formateurs freelance pour
assurer les cours en ligne. Ce modèle a des avantages évidents : un formateur humain
adapte son discours, répond aux questions, reformule une notion difficile, donne des
exemples et entretient une dynamique de groupe. Surtout, il incarne une responsabilité
pédagogique immédiate : si un élève ne comprend pas, il intervient sur-le-champ.

Mais ce modèle se heurte à un obstacle économique fort. Pour une petite structure, le coût
récurrent d'un formateur absorbe une part importante de la valeur produite par la
formation, et cette part croît avec le nombre d'heures dispensées. Dans un marché où les
marges doivent financer l'organisation, le suivi administratif, les outils et le
développement commercial, dépendre fortement d'un intervenant externe rend le modèle
difficile à industrialiser. Le problème initial peut donc se formuler simplement :
l'entreprise voulait développer une activité de formation externe, mais le modèle
classique du formateur en direct n'était compatible ni avec ses objectifs de coût, ni avec
son besoin de passage à l'échelle.

\subsection{Première réponse : séparer la production du cours de sa diffusion}

La première réponse a consisté à ne plus raisonner en cours en direct, mais en cours
préparé à l'avance. L'idée était de construire une plateforme capable d'héberger des
contenus audio et de les diffuser automatiquement le jour prévu de la formation. Plutôt
que de rémunérer un formateur pour animer plusieurs heures de cours, le travail était
scindé en deux étapes : la préparation du contenu, puis sa lecture sous forme d'audio.

Concrètement, les cours étaient rédigés pendant la semaine, éventuellement avec l'aide
d'outils d'intelligence artificielle, puis enregistrés et déposés dans la plateforme. Le jour
de la session, la diffusion se déclenchait selon un calendrier horodaté. Pour l'élève, le
cours se présentait comme une séance organisée dans le temps ; pour l'entreprise, la
production devenait plus facilement planifiable qu'un cours en direct. En répartissant
l'effort entre une personne qui écrit et une personne qui lit, le coût total devenait
nettement inférieur à celui d'un formateur freelance.

Cette organisation a validé plusieurs hypothèses importantes : les apprenants pouvaient
suivre une formation à partir de contenus audio diffusés à distance, l'entreprise pouvait
internaliser une partie de la production pédagogique, et la plateforme devenait l'outil
central d'orchestration de la formation, et non un simple espace de dépôt.

Elle ne résolvait toutefois qu'une partie du problème. La dépendance au formateur en
direct diminuait, mais la dépendance au travail humain demeurait. Les cours devaient
toujours être rédigés, relus, découpés, enregistrés, déposés et vérifiés. À chaque nouvelle
formation, voire à chaque mise à jour, une chaîne de production manuelle devait être
relancée. Le gain existait, mais il restait limité par le caractère artisanal du processus :
l'entreprise avait déplacé le coût humain plus qu'elle ne l'avait supprimé.

\subsection{De l'IA d'assistance à l'ambition d'une plateforme autonome}

L'arrivée des outils d'intelligence artificielle générative a déplacé la manière même de
poser le problème. Dans un premier temps, l'IA pouvait servir d'aide à la rédaction : un
collaborateur produisait une première version d'un cours, reformulait un passage,
structurait une séquence. L'humain gardait le contrôle et utilisait l'IA pour accélérer
certaines tâches.

Mais le contexte technologique a vite ouvert une hypothèse plus ambitieuse. Si un modèle
de langage peut rédiger un contenu, si un outil de synthèse vocale peut le lire, si un
système peut générer des supports visuels et si un assistant peut répondre aux questions
des élèves, alors il devient envisageable de concevoir une plateforme bien plus autonome.
Dans cette perspective, l'IA cesse d'être une aide ponctuelle pour devenir une composante
centrale de la chaîne de production et d'accompagnement.

Cet objectif doit être formulé sans détour : il s'agissait bien de remplacer une part
importante des interventions humaines initialement nécessaires à la formation. Dans une
petite entreprise, ce choix n'est pas seulement technologique, il est économique et
stratégique. L'enjeu est de rendre possible une activité qui serait trop coûteuse si elle
reposait entièrement sur des ressources humaines mobilisées en continu.

Remplacer une fonction humaine ne revient cependant pas à remplacer une personne par
un appel d'API. Un formateur ne se contente pas de lire un texte : il structure un cours,
choisit des exemples, respecte un programme, s'adapte au niveau des élèves, répond aux
questions, évalue la compréhension et maintient une cohérence entre les supports.
Automatiser cette fonction suppose donc d'automatiser une chaîne complète de tâches, et
non un geste isolé. C'est cette prise de conscience qui donne au sujet sa véritable
dimension d'ingénierie. La question n'est plus « quel outil d'IA utiliser ? », mais «
comment organiser une chaîne logicielle capable de produire, contrôler, diffuser et
exploiter des contenus pédagogiques générés par IA dans un contexte réel de formation ? »

\subsection{Les essais de synthèse vocale et leurs limites}

La première brique d'automatisation envisagée après la rédaction assistée a été la synthèse
vocale. L'idée était directe : si le cours est écrit à l'avance, une voix artificielle peut le
lire. Les outils récents produisent des voix bien plus naturelles que les anciens systèmes
de lecture automatique, ce qui rend la solution séduisante pour une plateforme audio.

Les premiers essais ont montré que le problème était plus subtil que prévu. Produire
plusieurs heures d'audio n'a rien d'une courte démonstration vocale. Les textes devaient
être découpés en portions compatibles avec l'outil, envoyés successivement, puis les
fragments audio obtenus devaient être réassemblés pour reconstituer un cours continu. Ce
procédé cumulait trois difficultés. Le temps de manipulation, d'abord : même générée par
IA, la voix exigeait un découpage et un assemblage largement manuels. Le coût, ensuite :
les meilleurs outils deviennent onéreux dès qu'on produit des journées complètes de
formation. La qualité pédagogique, enfin : une voix naturelle ne suffit pas si le texte lu
n'est pas pensé pour l'oral, s'il est répétitif ou s'il ne respecte pas le rythme d'un cours.

Ces essais ont clarifié un point essentiel : la synthèse vocale n'automatise que la lecture,
pas la conception du cours. Or un audio n'est exploitable que si le texte d'origine est clair,
progressif, calibré en durée, fidèle au programme et rédigé dans un style oral. La synthèse
vocale a donc déplacé le centre du problème vers l'amont. La vraie question n'était plus de
savoir comment générer une voix, mais comment générer un script pédagogique
suffisamment fiable pour être confié à une voix automatique et diffusé à de vrais élèves.

\subsection{Le cours préenregistré ne remplace pas le formateur}

Même avec un texte et une voix satisfaisants, le cours audio préenregistré bute sur une
limite structurelle : il ne peut pas interagir. Dans un cours en direct, l'élève pose une
question, demande une reformulation, signale une incompréhension, obtient un exemple
supplémentaire. Cette interaction fait partie intégrante de la valeur d'un formateur.

Dans un modèle purement audio, elle disparaît. Les élèves écoutent, mais ne disposent
d'aucun interlocuteur pour répondre à leurs questions. La limite est d'autant plus sensible
dans une formation professionnalisante, où les profils sont variés et les besoins de
clarification fréquents : un passage limpide pour l'un peut rester obscur pour l'autre.

Cette difficulté révèle une seconde dimension du problème. Il ne suffit pas d'automatiser
la production des cours ; il faut aussi automatiser une partie de l'accompagnement. Sans
cela, la plateforme reste un simple diffuseur d'audios, capable de réduire des coûts mais
incapable de remplacer les fonctions d'échange et de clarification d'un formateur. C'est
précisément le rôle d'un assistant conversationnel fondé sur la récupération d'information
(RAG) : non pas discuter librement comme un chatbot généraliste, mais répondre aux
questions des élèves à partir du contenu réel de la formation. La distinction est
fondamentale. Un assistant pédagogique ne doit pas produire une réponse simplement
plausible ; il doit s'appuyer sur les cours et les supports disponibles. Il doit être
contextualisé.

\subsection{Une classe de problèmes au-delà d'Audit Énergie}

Le cas d'Audit Énergie est singulier par son histoire et ses contraintes, mais le problème
qu'il pose dépasse largement cette entreprise. De nombreuses petites structures connaissent
aujourd'hui la même tension : elles veulent proposer des services plus riches, plus
personnalisés et plus accessibles, sans disposer des moyens humains pour les produire
manuellement. L'intelligence artificielle générative apparaît alors comme une issue : elle
sait rédiger, reformuler, résumer, produire de l'audio, générer des supports ou répondre à
des questions.

Mais son intégration dans un contexte professionnel ne peut reposer sur des usages
ponctuels. Un collaborateur qui demande à un modèle de rédiger un cours obtient parfois
un résultat utile, sans qu'il soit pour autant traçable, reproductible, contrôlé ni
directement exploitable en production. La classe de problèmes étudiée peut donc se
formuler ainsi : comment transformer des capacités d'IA générative, le plus souvent
utilisées de façon ponctuelle, en une chaîne logicielle fiable, auditable et intégrée à un
processus métier réel ?

Dans le cas de la formation, cette question prend une forme particulièrement exigeante. Il
ne s'agit pas seulement de produire du texte, mais un objet pédagogique complet : un cours
qui respecte un programme, s'adapte au niveau des apprenants, reste compréhensible,
tient dans une durée, peut être transformé en audio, donne lieu à des supports visuels,
permet une évaluation et s'inscrit dans un parcours plus large. L'IA générative peut ici
donner une illusion de simplicité : un premier résultat convaincant masque les difficultés
réelles de répétabilité, de contrôle qualité, de coût et d'évaluation. Pour un prototype, il
suffit qu'un texte semble correct ; pour une plateforme utilisée avec de vrais élèves, il faut
pouvoir expliquer comment le contenu a été produit, quelles règles ont été appliquées et
comment sa qualité est mesurée.

\subsection{Les exigences qui encadrent le problème}

Trois exigences encadrent ce problème et orientent toute la suite du mémoire.

La première est la \emph{qualité pédagogique}. Dans une formation professionnelle, un
contenu n'est pas jugé sur sa fluidité ou son apparence, mais sur sa capacité à faire
progresser l'apprenant vers un objectif précis. Or les modèles de langage produisent des
textes très convaincants, qui peuvent pourtant contenir des imprécisions, des répétitions,
des exemples mal choisis ou des informations non vérifiées. Le risque est concret : un
élève peut retenir une explication inexacte ou perdre du temps sur un contenu hors
programme. Cet enjeu se double d'un problème de \emph{densité}. Lorsqu'on demande à
la chaîne de produire un long volume de cours à partir d'une source documentaire limitée,
le modèle, contraint de remplir le temps, tend à broder, à répéter et à digresser. La qualité
ne peut donc pas être traitée par une simple relecture finale : elle doit être intégrée à la
chaîne de production elle-même.

La deuxième exigence est la \emph{traçabilité}. Lorsqu'un humain produit un cours, il
peut expliquer ses choix et justifier sa progression. Lorsqu'une IA génère un contenu, cette
justification n'existe pas spontanément. Si la plateforme se contente de stocker le texte
final, il devient très difficile de savoir d'où vient une erreur : du programme de départ, du
plan généré, de la rédaction d'une section, d'une étape de réécriture ou de la calibration
audio. Sans artefacts intermédiaires, ni la correction, ni la confiance, ni l'amélioration
continue ne sont possibles. La traçabilité n'est pas un confort : c'est une condition de
l'évaluation.

La troisième exigence est d'ordre \emph{économique et organisationnel}. L'objectif n'est
pas de construire une plateforme techniquement impressionnante, mais de rendre viable un
modèle de formation compatible avec les moyens d'une petite entreprise. L'automatisation
ne supprime pas tous les coûts : les services d'IA ont un prix, les erreurs doivent être
contrôlées, les prompts maintenus, la plateforme supervisée et les contenus validés.
Évaluer le projet ne reviendra donc pas à comparer un formateur humain à un modèle
d'IA, mais à comparer deux organisations : l'une centrée sur l'intervention humaine
continue, l'autre sur une plateforme automatisée supervisée.

\subsection{Formulation de la problématique}

Le problème posé par ce mémoire se lit à deux niveaux. Au niveau d'Audit Énergie,
l'entreprise veut développer une formation externe avec des moyens humains limités : le
formateur en direct coûte trop cher, le cours préenregistré appauvrit l'interaction
pédagogique. Elle a donc besoin d'une plateforme capable d'automatiser la production, la
diffusion et une partie de l'accompagnement. À un niveau plus général, le même problème
concerne toutes les petites structures qui cherchent à industrialiser un service grâce à
l'IA : l'enjeu n'est pas tant d'utiliser ces outils que de les encadrer dans une architecture
permettant un usage fiable en production.

La problématique peut alors être formulée ainsi :

\begin{quote}
\emph{Comment concevoir une plateforme de formation autonome alimentée par
l'intelligence artificielle, capable de remplacer une part significative des interventions
humaines nécessaires à la production, à la diffusion et à l'accompagnement des cours,
tout en garantissant la qualité pédagogique, la traçabilité des contenus et la pertinence
des réponses apportées aux apprenants ?}
\end{quote}

Cette formulation met en évidence la tension centrale du projet. D'un côté, l'entreprise
cherche à automatiser fortement son modèle pour le rendre économiquement viable. De
l'autre, cette automatisation ne peut se faire au prix d'une perte de qualité, d'une absence
de contrôle ou d'une expérience dégradée pour les élèves. Le travail d'ingénierie consiste
donc à concilier autonomie et maîtrise.

\subsection{Objectifs de la solution attendue}

La problématique se décline en objectifs concrets, qui structureront l'étude de la solution :

\begin{itemize}
\item \textbf{Automatiser la génération des cours} à partir d'un cadre pédagogique donné,
en respectant une progression et un niveau attendus. Cette génération ne doit pas être une
rédaction libre, mais un processus structuré et contrôlable.
\item \textbf{Automatiser la transformation audio}, en calibrant les contenus pour la
lecture orale et en évitant tout montage manuel lourd, là où chaque fichier exigeait
auparavant une manipulation individuelle.
\item \textbf{Produire des supports complémentaires} (slides, exercices) cohérents avec le
cours, et non ajoutés après coup comme des éléments indépendants.
\item \textbf{Accompagner les apprenants} via un assistant conversationnel qui répond à
partir du contenu réel de la formation, plutôt que d'une connaissance générale non
contrôlée.
\item \textbf{Rendre le système auditable}, en associant à chaque contenu ses traces :
plan, sections, corrections, sources et indicateurs.
\item \textbf{Définir des indicateurs d'évaluation} : temps de production, nombre
d'interventions humaines, coût, respect des durées, qualité des contenus, résultats aux
exercices et pertinence des réponses de l'assistant.
\end{itemize}

Ces objectifs ne sont pas seulement des résultats à présenter en fin de mémoire : ils
orientent la conception. Pour mesurer le respect d'un plan, il faut un plan explicite ; pour
mesurer les corrections, il faut conserver les versions intermédiaires ; pour mesurer la
pertinence du RAG, il faut pouvoir identifier les sources utilisées. L'évaluation
influence donc directement l'architecture.

\subsection{Motivation, périmètre et transition}

Ce sujet présente un intérêt particulier pour un mémoire d'ingénieur, car il se situe à la
frontière d'un besoin métier concret et de problématiques techniques actuelles. Il ne s'agit
pas d'étudier l'intelligence artificielle de manière abstraite, mais de l'intégrer dans un
système réellement utilisé, avec de vrais élèves, des contraintes de coût et des exigences de
qualité. Sa valeur tient aussi à sa progression : la plateforme n'est pas née sous sa forme
actuelle, elle a évolué par étapes — diffusion d'audios, préparation manuelle, essais de
synthèse vocale, puis structuration d'une pipeline de génération et ajout d'un assistant
RAG. Chaque étape a fait surgir une nouvelle limite et obligé à reformuler le problème.
Prendre du recul sur cette trajectoire permet d'expliquer non seulement ce qui a été
développé, mais pourquoi ces choix ont été faits et quels compromis ont été acceptés.

Le périmètre du mémoire reste délimité. Il porte sur la conception d'une plateforme
autonome pour la production, la diffusion et l'accompagnement de cours à distance ; les
aspects administratifs, commerciaux ou réglementaires de l'organisme de formation ne sont
abordés que lorsqu'ils éclairent les contraintes du projet. Le travail ne prétend pas non
plus démontrer que l'IA peut remplacer toute forme d'enseignement humain : la
supervision humaine demeure nécessaire, mais elle change de rôle, davantage tournée vers
la validation, le pilotage et le traitement des cas qui dépassent le cadre prévu.

Cette position du problème fait apparaître plusieurs familles de solutions à examiner :
les plateformes de formation en ligne, qui diffusent et suivent mais ne produisent pas les
cours ; les outils de synthèse vocale, qui automatisent la lecture mais pas la conception ;
les modèles de langage, qui génèrent du texte mais doivent être encadrés ; les systèmes
RAG, qui contextualisent les réponses mais dépendent de la qualité de la recherche
documentaire. L'état de l'art, présenté dans la partie suivante, analysera ces approches
pour mieux situer, ensuite, la solution développée et en montrer l'originalité.

% ====================== PARTIE 2 — ÉTAT DE L'ART ======================
\clearpage
\section{État de l'art}

\subsection{Positionnement de l'état de l'art}

Le problème posé par ce mémoire n'est pas seulement celui de l'utilisation d'un chatbot
ou d'un générateur de texte. Il concerne la conception d'une plateforme de formation
autonome alimentée par l'intelligence artificielle. Deux dimensions doivent donc être
étudiées conjointement. La première concerne l'assistant pédagogique : comment permettre
à un modèle de langage de répondre aux questions des apprenants à partir des contenus
réels de la formation ? La seconde concerne la production de contenus : comment générer
des cours longs, structurés, oralisables et contrôlables sans se limiter à une simple sortie
textuelle produite par un prompt ?

L'état de l'art est volontairement centré sur les techniques plutôt que sur un catalogue
d'outils commerciaux. Les plateformes d'apprentissage en ligne permettent déjà de
diffuser des cours, de suivre des élèves ou de proposer des évaluations, mais elles ne
répondent pas directement à la question de recherche de ce mémoire : rendre une chaîne
de formation suffisamment autonome tout en conservant une qualité pédagogique, une
traçabilité et une capacité d'évaluation. Les approches étudiées ici sont donc celles qui
permettent d'adapter ou d'encadrer les modèles de langage dans un contexte métier.

La progression adoptée est la suivante. On étudie d'abord les approches permettant de
donner un contexte à un modèle de langage : prompt engineering, fine-tuning et RAG. On
analyse ensuite les variantes du RAG, notamment la recherche vectorielle, la recherche
hybride, le reclassement sémantique, les métadonnées et l'évaluation. Enfin, on étudie
les méthodes de génération contrôlée de contenus longs : décomposition de tâche,
génération structurée, schémas intermédiaires, auto-évaluation et validation humaine. Ce
cheminement prépare l'étude originale, sans encore détailler la solution développée.

\subsection{Modèles de langage et adaptation au contexte}

Les grands modèles de langage ont montré qu'ils peuvent réaliser de nombreuses tâches
à partir d'instructions en langage naturel et de quelques exemples fournis dans le
prompt. Les travaux sur le few-shot learning ont mis en évidence cette capacité à adapter
le comportement du modèle sans modifier ses poids \cite{brown2020language}. Pour une
plateforme de formation, cette propriété est importante : un modèle généraliste peut être
orienté vers un rôle d'assistant pédagogique, de rédacteur de cours ou de correcteur de
contenu simplement à partir d'instructions.

Cette flexibilité explique l'intérêt du prompt engineering. Dans sa forme la plus simple,
il consiste à fournir au modèle une consigne, un rôle, des contraintes de style, un format
de sortie et parfois des exemples. Dans le cas d'un chatbot pédagogique, le prompt peut
indiquer que le modèle doit répondre clairement, ne pas inventer d'information et rester
dans le périmètre du cours. Dans le cas d'une génération de contenu, il peut préciser le
niveau de l'apprenant, le ton attendu, la structure du cours ou les règles de conformité.

Le prompt engineering présente trois avantages. Il est rapide à mettre en place, peu
coûteux à modifier et compatible avec des modèles propriétaires accessibles par API. Il
est donc très utile en phase de prototypage. Cependant, il a aussi des limites fortes. Le
contexte disponible dans le prompt reste limité ; les consignes peuvent entrer en
concurrence ; et le modèle peut suivre certaines instructions tout en en négligeant
d'autres. Lorsque le volume documentaire augmente ou lorsque les règles métier deviennent
nombreuses, le prompt devient difficile à maintenir.

Surtout, le prompt ne constitue pas une mémoire documentaire fiable. Il permet de
fournir ponctuellement des informations, mais il ne garantit pas que le modèle réponde
à partir d'une source précise. Dans une formation, cette limite est critique : l'assistant
doit répondre à partir des contenus réellement enseignés, pas à partir d'une connaissance
générale ou approximative. De même, un générateur de cours doit respecter un programme
et ne pas produire une progression plausible mais non conforme.

\subsection{Fine-tuning et spécialisation du modèle}

Une autre approche consiste à adapter le modèle par fine-tuning. Au lieu de fournir toutes
les règles dans le prompt, on entraîne le modèle sur des exemples représentatifs de la
tâche attendue. Cette approche peut améliorer le style, le format de sortie ou la
répétabilité d'un comportement. Elle est particulièrement utile lorsque l'on souhaite
spécialiser un modèle sur un type de réponse ou un domaine récurrent.

Le fine-tuning classique peut cependant être coûteux, car il nécessite de modifier un
grand nombre de paramètres. Des méthodes plus légères, comme LoRA, cherchent à
réduire ce coût en gelant les poids du modèle et en ajoutant des matrices de bas rang
entraînables \cite{hu2021lora}. Ces approches rendent la spécialisation plus accessible,
mais elles ne suppriment pas les contraintes fondamentales du fine-tuning : il faut
constituer un jeu de données de qualité, entraîner ou adapter un modèle, évaluer le
résultat et maintenir cette adaptation dans le temps.

Dans le contexte d'un assistant pédagogique, le fine-tuning n'est pas toujours la meilleure
réponse au problème de connaissance. Il permet d'apprendre un style ou une procédure,
mais il est moins adapté pour intégrer des documents qui changent souvent. Si un cours,
un référentiel ou un support est modifié, il n'est pas souhaitable de réentraîner le modèle
à chaque mise à jour. En outre, le fine-tuning ne garantit pas que la réponse produite soit
fidèle à un document précis. Un modèle fine-tuné peut continuer à halluciner ou à
mélanger des informations apprises avec des informations plausibles.

Pour une plateforme de formation, le fine-tuning est donc pertinent pour stabiliser un
comportement, mais insuffisant comme mécanisme principal de contextualisation
documentaire. Il peut apprendre au modèle à répondre comme un assistant pédagogique,
mais il ne suffit pas à lui donner une mémoire fiable, actualisable et traçable des cours.

\subsection{RAG : ancrer la génération dans des documents}

Le Retrieval-Augmented Generation, ou RAG, répond à cette limite en séparant deux
fonctions : la récupération d'information et la génération de réponse. Le principe est de
ne pas modifier directement le modèle, mais de rechercher au moment de la question des
passages pertinents dans une base documentaire, puis de fournir ces passages au modèle
comme contexte. Le modèle génère alors une réponse à partir des documents récupérés.
Cette approche a été formalisée dans les travaux de Lewis et al. sur les tâches de NLP
intensives en connaissances \cite{lewis2020retrieval}.

Le RAG présente plusieurs avantages pour un assistant pédagogique. D'abord, les contenus
peuvent être mis à jour sans réentraîner le modèle. Il suffit de réindexer les documents.
Ensuite, la réponse peut être rattachée à des passages sources, ce qui améliore la
traçabilité. Enfin, le modèle est moins dépendant de sa mémoire paramétrique : il peut
répondre à partir du cours réellement fourni, même si ce contenu est spécifique à une
formation ou à une promotion.

Cette architecture est particulièrement adaptée au contexte de ce mémoire. Les élèves ne
doivent pas recevoir une réponse générique sur un métier ou une notion ; ils doivent
recevoir une réponse cohérente avec le cours diffusé, le référentiel utilisé et les supports
mis à disposition. Le RAG permet donc de transformer un modèle généraliste en assistant
contextualisé.

Cependant, le RAG ne supprime pas tous les risques. Il déplace une partie du problème
vers la qualité de la récupération. Si les passages récupérés sont incomplets, trop larges,
hors sujet ou mal classés, le modèle peut produire une réponse faible. De plus, même
avec un bon contexte, le modèle peut reformuler de manière imprécise ou ajouter une
information non supportée par les documents. Les travaux sur les hallucinations montrent
que les modèles de langage peuvent générer des informations plausibles mais non fondées,
ce qui impose des mécanismes d'évaluation et de contrôle \cite{ji2023hallucination}.

\subsection{Techniques de recherche dans les architectures RAG}

Un système RAG n'est pas une technique unique. Il regroupe plusieurs choix d'architecture
qui influencent directement la qualité des réponses. Le premier choix concerne le
découpage documentaire. Les documents sont rarement injectés entiers dans le modèle :
ils sont découpés en segments, ou chunks, qui seront indexés puis récupérés. Des chunks
trop petits peuvent perdre le contexte nécessaire à la compréhension ; des chunks trop
grands peuvent diluer l'information pertinente. Le choix de la taille, du chevauchement
et des métadonnées est donc déterminant.

Le deuxième choix concerne la représentation des documents. Une recherche par mots-clés
fonctionne bien lorsque la question reprend le vocabulaire exact du document, mais elle
peut échouer lorsque l'élève reformule avec d'autres termes. La recherche vectorielle
répond à cette limite en représentant les questions et les passages dans un espace
sémantique. Elle permet de retrouver des passages proches par le sens, même si les mots
ne sont pas identiques.

La recherche vectorielle a elle aussi des limites. Elle peut retrouver des passages
sémantiquement proches mais pas exactement pertinents, ou négliger des mots clés
importants comme une référence, un code, une notion réglementaire ou un terme métier.
C'est pourquoi de nombreuses architectures combinent recherche lexicale et recherche
vectorielle. On parle alors de recherche hybride. Les travaux récents sur le RAG décrivent
cette évolution d'un RAG naïf vers des architectures plus avancées ou modulaires
\cite{gao2023retrieval}.

Le troisième choix concerne le reclassement des résultats. Après une première recherche,
un système peut reranker les passages candidats pour privilégier ceux qui correspondent
le mieux à l'intention de la question. Ce reclassement peut s'appuyer sur des modèles de
langage ou des systèmes spécialisés. Dans un contexte pédagogique, il est utile car les
questions des élèves sont souvent courtes, imprécises ou formulées avec des termes non
experts.

Le quatrième choix concerne les métadonnées. Un même corpus peut contenir plusieurs
formations, plusieurs jours, plusieurs modules ou plusieurs types de documents. Sans
filtrage, l'assistant risque de récupérer un passage provenant d'un mauvais cours. Les
métadonnées permettent de restreindre la recherche à un périmètre pertinent : formation,
jour, module, type de support ou version du document.

\subsection{Azure AI Search comme approche industrielle du RAG}

Plusieurs solutions peuvent être utilisées pour construire un RAG : base vectorielle locale,
index spécialisé, moteur de recherche hybride ou service cloud managé. Dans un prototype,
une bibliothèque locale comme FAISS peut suffire pour tester la recherche vectorielle.
Dans une plateforme destinée à être utilisée par de vrais élèves, d'autres critères
deviennent importants : indexation automatique, recherche hybride, filtrage, supervision,
sécurité, intégration avec le stockage documentaire et capacité à maintenir plusieurs
corpus.

Azure AI Search s'inscrit dans cette logique industrielle. La documentation Microsoft
présente Azure AI Search comme un service d'indexation et de recherche pouvant être
utilisé dans des architectures RAG avec des données textuelles et vectorielles
\cite{microsoft2026foundryrag}. Le service permet de combiner recherche par mots-clés
et recherche vectorielle dans une requête hybride ; les résultats sont fusionnés avant
d'être exploités par l'application \cite{microsoft2026hybridsearch}. Il propose également
un semantic ranker, qui reclasse les résultats textuels ou hybrides à partir de modèles de
compréhension du langage \cite{microsoft2026semanticranker}.

Dans le contexte d'un chatbot pédagogique, ces fonctionnalités répondent à plusieurs
besoins. La recherche hybride permet de traiter à la fois les formulations approximatives
des élèves et les termes précis présents dans les supports. Le semantic ranking permet
d'améliorer la pertinence des passages présentés au modèle. Les index et métadonnées
permettent de séparer les formations ou les plateformes. Enfin, l'intégration avec Azure
OpenAI simplifie la construction d'une chaîne de question-réponse où les documents
indexés sont fournis comme source au modèle.

Cette approche a toutefois des limites. Un service managé réduit la complexité
d'infrastructure, mais il ne dispense pas de concevoir correctement le corpus. La qualité
du RAG dépend encore du découpage, des métadonnées, des documents indexés, de la
formulation des requêtes et du prompt final. Le semantic ranker reclasse des résultats
déjà récupérés ; il ne peut pas compenser entièrement une mauvaise indexation ou un
corpus incomplet. Le choix d'Azure AI Search doit donc être compris comme un choix
d'architecture robuste, pas comme une garantie automatique de qualité.

\subsection{Évaluation des systèmes RAG}

L'évaluation d'un RAG doit distinguer la qualité de la récupération et la qualité de la
réponse. Une réponse peut être bien formulée mais fondée sur de mauvais passages ; à
l'inverse, les bons passages peuvent être récupérés mais mal exploités par le modèle.
Plusieurs cadres d'évaluation récents proposent donc de mesurer séparément la pertinence
du contexte, la fidélité de la réponse aux sources et la pertinence de la réponse finale.
C'est le cas de RAGAS \cite{es2023ragas} et d'ARES \cite{saadfalcon2023ares}.

Dans un contexte pédagogique, les métriques utiles peuvent être regroupées en quatre
familles. La première concerne la récupération : précision@k, rappel@k, rang moyen du
bon passage, taux de questions pour lesquelles au moins une source pertinente est
retrouvée. La deuxième concerne la fidélité : la réponse est-elle supportée par les passages
fournis ? Ajoute-t-elle une information non présente dans les documents ? La troisième
concerne l'utilité pédagogique : la réponse est-elle claire, adaptée au niveau de l'élève et
suffisamment concise ? La quatrième concerne l'expérience réelle : les élèves reposent-ils
la même question, utilisent-ils l'assistant et jugent-ils les réponses utiles ?

La citation des sources est également un enjeu. Les travaux sur la vérifiabilité dans les
moteurs génératifs montrent que les systèmes doivent citer de manière complète et
précise les éléments qui soutiennent leurs réponses \cite{liu2023evaluating}. Pour une
formation, cette exigence est moins formelle que dans une recherche documentaire
académique, mais elle reste importante : un responsable pédagogique doit pouvoir
vérifier d'où vient une réponse.

Ces considérations conduisent à une conclusion claire : un RAG ne doit pas être évalué
uniquement par impression subjective. Il faut construire un jeu de questions représentatif,
identifier les passages attendus, comparer plusieurs configurations de recherche et
mesurer la fidélité des réponses. Sans cette évaluation, il est impossible de justifier
scientifiquement le choix d'une architecture RAG plutôt qu'une autre.

\subsection{Génération de contenus longs par prompt unique}

La deuxième dimension du mémoire concerne la génération automatique des contenus de
formation. La première approche possible consiste à demander directement à un modèle
de langage de produire un cours complet à partir d'un programme, d'un référentiel et de
quelques consignes. Cette approche est séduisante parce qu'elle exploite la capacité des
modèles à rédiger des textes cohérents et à suivre un format général.

Dans un contexte simple, la génération directe peut donner des résultats acceptables. Elle
est utile pour produire un premier brouillon, un résumé, un plan de cours ou une
explication courte. Elle montre rapidement que l'IA peut réduire le temps de rédaction.
Cependant, elle devient fragile lorsque la sortie attendue est longue, structurée,
oralisable et destinée à être utilisée sans réécriture complète.

Plusieurs problèmes apparaissent alors. Le modèle peut répéter des idées, perdre le fil de
la progression, produire des transitions artificielles, dépasser la longueur attendue ou
oublier certaines contraintes. Il peut aussi mélanger le langage interne du système avec
le langage destiné à l'apprenant, par exemple en verbalisant des contraintes techniques
qui ne devraient pas apparaître dans le cours. Enfin, si le résultat final est mauvais, il est
difficile de savoir quelle partie du prompt ou de la génération est responsable.

Ces limites ne signifient pas que le prompt engineering est inutile. Elles montrent plutôt
que le prompt unique atteint rapidement ses limites lorsqu'il doit porter toutes les
responsabilités : structure pédagogique, style oral, durée, conformité, exemples,
supports, exercices et qualité. Un prompt trop chargé devient un artefact fragile, difficile
à faire évoluer et difficile à évaluer.

\subsection{Décomposition de tâches et génération structurée}

Pour dépasser ces limites, de nombreux travaux sur les modèles de langage insistent sur
la décomposition des tâches complexes. Le chain-of-thought a montré que faire produire
des étapes intermédiaires peut améliorer certaines tâches de raisonnement
\cite{wei2022chain}. Le least-to-most prompting va plus loin en demandant au modèle de
décomposer un problème en sous-problèmes plus simples avant de les résoudre
\cite{zhou2022least}. Les approches de type Tree of Thoughts proposent d'explorer
plusieurs pistes intermédiaires avant de choisir une solution \cite{yao2023tree}.

Ces travaux portent principalement sur le raisonnement, mais leur logique est applicable
à la génération de contenus pédagogiques longs. Un cours complet est une tâche complexe
qui peut être décomposée : analyser le programme, construire un plan, détailler les
sections, rédiger les transitions, adapter le texte à l'oral, générer des supports, produire
des questions, puis vérifier la conformité. La qualité peut alors être contrôlée à chaque
étape plutôt qu'uniquement à la fin.

La génération structurée repose sur la même idée. Au lieu de produire directement un
texte final, le modèle produit d'abord une représentation intermédiaire : plan, JSON,
liste d'objectifs, sections, contraintes de durée ou points clés. Cette représentation sert
de contrat entre les étapes. Elle permet d'auditer la structure avant de rédiger, de répartir
les budgets, de générer section par section et de rattacher les supports à des intentions
pédagogiques explicites.

Cette approche présente des avantages importants. Elle réduit la charge imposée au
modèle à chaque appel, rend les erreurs plus localisables et facilite l'intégration logicielle.
Elle introduit aussi une forme de traçabilité : on peut comparer le plan prévu, le texte
généré et les corrections appliquées. En revanche, elle augmente la complexité de
l'orchestration. Il faut gérer plusieurs appels au modèle, valider les formats, traiter les
erreurs JSON, conserver les artefacts et maîtriser le coût.

\subsection{Auto-évaluation, reviews automatiques et validation humaine}

Une autre famille d'approches consiste à utiliser les modèles de langage non seulement
pour générer, mais aussi pour critiquer ou améliorer leurs propres sorties. Des travaux
comme Reflexion explorent l'idée d'agents capables d'utiliser un retour verbal pour
améliorer leurs actions \cite{shinn2023reflexion}. Dans les applications de génération de
contenu, cette idée se traduit souvent par des passes de review : un modèle relit le texte,
identifie des problèmes, propose des corrections ou vérifie le respect de règles.

Les reviews automatiques sont intéressantes pour une plateforme de formation, car elles
permettent de contrôler plusieurs dimensions : cohérence avec le plan, ton pédagogique,
adaptation à l'oral, conformité éthique, absence de formulations sensibles, respect de la
durée ou qualité des QCM. Elles peuvent réduire la charge de validation humaine en
signalant les passages problématiques.

Elles ne doivent toutefois pas être considérées comme infaillibles. Un modèle peut ne pas
détecter une erreur, proposer une correction qui dégrade le contenu ou valider une
réponse insuffisamment sourcée. L'évaluation par LLM-as-a-judge peut être utile, mais
elle possède des biais et doit être encadrée par des critères explicites \cite{zheng2023judging}.
Dans un système rigoureux, les reviews automatiques doivent donc être combinées avec
des règles déterministes, des métriques et, lorsque le risque le justifie, une validation
humaine.

La validation humaine garde une place importante dans les systèmes pédagogiques. Elle
peut intervenir à plusieurs niveaux : validation du programme, relecture d'échantillons,
contrôle des sorties à risque, analyse des retours élèves et amélioration des règles. Le
but n'est pas nécessairement de relire chaque phrase produite, mais de déplacer
l'intervention humaine vers les points où elle apporte le plus de valeur.

\subsection{Évaluation des contenus pédagogiques générés}

L'évaluation des contenus générés est plus difficile que l'évaluation d'une tâche fermée.
Un cours n'a pas une seule bonne réponse. Il faut donc combiner plusieurs indicateurs.
Le premier est la fidélité au programme : les compétences prévues sont-elles couvertes ?
Les notions importantes sont-elles présentes ? Le cours évite-t-il les informations hors
périmètre ?

Le deuxième indicateur est la cohérence pédagogique. Un cours doit suivre une progression
compréhensible, éviter les répétitions inutiles, introduire les notions avant de les utiliser
et proposer des exemples adaptés. Cette dimension peut être évaluée par relecture humaine,
par grille critériée ou par comparaison entre le plan et le texte généré.

Le troisième indicateur est l'adaptation à l'oral. Un texte destiné à la synthèse vocale ne
doit pas être écrit comme un support PDF. Il doit contenir des phrases prononçables, des
transitions, un rythme d'écoute et des formulations naturelles. La durée est également
un critère : si le texte dépasse le créneau audio prévu, la formation devient difficile à
diffuser.

Le quatrième indicateur concerne les supports associés. Les slides doivent correspondre
aux moments pédagogiques importants, et les QCM doivent vérifier des notions réellement
enseignées. Un QCM généré indépendamment du cours peut donner une illusion
d'évaluation sans mesurer la compréhension du contenu diffusé.

Enfin, il faut mesurer l'impact opérationnel. Une pipeline de génération n'a d'intérêt que
si elle réduit effectivement le temps de production, les manipulations manuelles ou le
coût, tout en maintenant une qualité acceptable. Les métriques doivent donc combiner
qualité pédagogique et efficacité industrielle.

\subsection{Synthèse critique}

Les approches étudiées montrent que les modèles de langage ne suffisent pas, à eux
seuls, à construire une plateforme de formation autonome. Le prompt engineering est
utile pour orienter le comportement du modèle, mais il devient fragile lorsque les règles
et les documents se multiplient. Le fine-tuning permet d'adapter un style ou une tâche,
mais il n'est pas idéal pour injecter des connaissances documentaires souvent mises à
jour. Le RAG apporte une solution plus adaptée pour l'assistant pédagogique, car il
ancre les réponses dans les contenus de formation et permet une mise à jour des sources.

De la même manière, la génération directe d'un cours par prompt unique est insuffisante
pour produire des contenus longs, oralisables et contrôlables. Les approches de
décomposition, de génération structurée et de review automatique offrent une meilleure
base pour construire une pipeline fiable. Elles transforment la génération en processus,
avec des étapes, des artefacts et des points de contrôle.

Le point commun entre ces deux axes est la nécessité d'encadrer l'IA. Dans un assistant
RAG, l'encadrement passe par la récupération, les sources et l'évaluation de la fidélité.
Dans une pipeline de génération, il passe par le plan, les formats structurés, les reviews
et les métriques. L'originalité de l'approche étudiée dans la suite du mémoire devra donc
être évaluée à partir de cette double exigence : contextualiser les réponses du chatbot et
structurer la génération des contenus.

\subsection{Tableau récapitulatif des approches}

\begin{table}[h]
\centering
\small
\begin{tabular}{|p{3cm}|p{3.4cm}|p{4.2cm}|p{3.4cm}|}
\hline
\textbf{Approche} & \textbf{Apport principal} & \textbf{Limites} & \textbf{Intérêt pour le projet} \\
\hline
Prompt engineering &
Orientation rapide du modèle par consignes, rôles, exemples et contraintes. &
Contexte limité, maintenance difficile, faible traçabilité, risque d'oubli de consignes. &
Utile pour le style et les règles locales, insuffisant seul pour une plateforme autonome. \\
\hline
Fine-tuning &
Adaptation du comportement ou du style du modèle sur des exemples. &
Coût de constitution des données, maintenance lourde, pas idéal pour documents changeants. &
Pertinent pour stabiliser un comportement, moins adapté à la connaissance pédagogique actualisable. \\
\hline
RAG simple &
Recherche de passages pertinents puis génération ancrée dans les documents. &
Dépend du chunking, de la qualité du corpus et de la pertinence de la récupération. &
Base pertinente pour répondre aux élèves à partir des contenus de cours. \\
\hline
RAG avancé &
Recherche hybride, reranking, métadonnées, citations et évaluation. &
Architecture plus complexe, coûts supplémentaires, besoin de jeux de tests. &
Approche la plus adaptée pour un assistant pédagogique traçable et maintenable. \\
\hline
Azure AI Search &
Service managé combinant indexation, recherche vectorielle, recherche hybride et semantic ranker. &
Ne compense pas un mauvais corpus ; dépend de la configuration, des filtres et de l'évaluation. &
Choix cohérent pour industrialiser le RAG dans une plateforme utilisée par de vrais élèves. \\
\hline
Génération directe de cours &
Production rapide de brouillons ou de contenus courts. &
Dérive sur contenus longs, répétitions, contrôle faible de la durée et de la structure. &
Approche utile au prototypage, mais insuffisante comme chaîne de production. \\
\hline
Génération structurée &
Plan, sections, schémas et artefacts intermédiaires. &
Orchestration plus complexe, validation JSON et coût d'appels multiples. &
Approche adaptée à la production de cours longs, audios, slides et QCM. \\
\hline
Reviews automatiques &
Contrôle du style, de la conformité, de l'adhérence au plan et de la qualité. &
Risque de faux positifs ou faux négatifs ; nécessite métriques et supervision. &
Permet de réduire la relecture humaine tout en rendant la qualité plus observable. \\
\hline
Human-in-the-loop &
Validation humaine sur les points critiques. &
Réintroduit du coût humain et limite l'autonomie complète. &
Nécessaire pour superviser le système, surtout en production pédagogique. \\
\hline
\end{tabular}
\caption{Synthèse des approches existantes et de leur intérêt pour une plateforme de formation autonome}
\end{table}

\subsection{Comparaisons et évaluations à réaliser}

Pour rendre le mémoire scientifiquement rigoureux, plusieurs comparaisons doivent être
réalisées ou, à défaut, présentées comme des évaluations à consolider. La première
comparaison concerne le chatbot : prompt seul contre RAG. Il s'agit de construire un
jeu de questions d'élèves, de définir les passages sources attendus, puis de comparer la
fidélité, la pertinence et le taux d'hallucination des réponses.

La deuxième comparaison concerne les variantes de RAG. Il serait pertinent de comparer
une recherche vectorielle simple, une recherche hybride et une recherche hybride avec
reclassement sémantique lorsque la configuration le permet. Les métriques attendues
sont la précision@k, le rappel@k, le rang du premier passage pertinent et la fidélité de
la réponse finale.

La troisième comparaison concerne la génération de contenus. Une génération directe par
prompt unique peut être comparée à une génération structurée par plan et sections. Les
critères peuvent porter sur le respect du programme, la répétition, la longueur, la
cohérence de la progression, la qualité orale et le nombre de corrections nécessaires.

La quatrième comparaison concerne les reviews automatiques. Il faut mesurer ce qu'elles
détectent réellement : nombre de problèmes signalés, taux de corrections acceptées,
effets indésirables sur le texte et réduction de la charge de relecture humaine. Cette
évaluation est importante pour éviter de présenter les reviews comme une garantie de
qualité sans preuve.

La cinquième comparaison concerne l'impact opérationnel. Il faut comparer le temps de
production et le nombre d'interventions humaines avant et après automatisation. Même si
les chiffres restent partiels, cette mesure est indispensable pour relier la solution à la
problématique économique de l'entreprise.

% ====================== PARTIE 3 — ÉTUDE ORIGINALE ======================
\clearpage
\section{Étude originale}

\subsection{Objectif de l'approche proposée}

L'état de l'art a montré que deux problèmes doivent être traités séparément pour
construire une plateforme de formation autonome alimentée par l'intelligence
artificielle. Le premier est un problème de contextualisation : un assistant pédagogique
ne doit pas répondre comme un chatbot généraliste, mais à partir des contenus réellement
disponibles dans la formation. Le second est un problème de génération contrôlée : un
modèle de langage peut produire du texte, mais la production d'une journée complète de
formation audio, accompagnée de supports et d'exercices, nécessite une chaîne plus
structurée qu'un simple prompt.

L'approche proposée dans ce mémoire consiste donc à combiner deux architectures
complémentaires. D'une part, un assistant pédagogique fondé sur une logique RAG, afin
de répondre aux questions des élèves à partir des documents indexés de la formation.
D'autre part, une pipeline de génération de contenus pédagogiques, destinée à transformer
un référentiel ou un programme de formation en contenus audio, supports visuels et
artefacts contrôlables.

L'originalité de l'approche ne réside pas dans l'utilisation isolée d'un modèle de langage
ou d'un service de synthèse vocale. Ces briques sont désormais disponibles dans de
nombreux environnements. Elle réside plutôt dans l'orchestration de ces briques dans une
plateforme utilisable par de vrais élèves, avec une attention portée à la traçabilité, à la
qualité pédagogique, au coût et à la capacité de réutiliser les modules générés. La solution
ne cherche pas seulement à produire une réponse ou un texte ; elle cherche à produire un
objet pédagogique exploitable.

L'hypothèse principale est la suivante : une plateforme de formation peut remplacer une
part significative des interventions humaines répétitives si l'IA est encadrée par une
architecture en étapes, des contrats de données explicites, des mécanismes de contrôle
qualité et des indicateurs d'évaluation. L'autonomie ne doit donc pas être comprise comme
une absence totale de supervision humaine, mais comme une réduction forte du travail
manuel récurrent grâce à une chaîne automatisée observable.

\subsection{Modèle fonctionnel de la plateforme}

La plateforme développée répond à quatre grands cas d'utilisation. Le premier concerne
l'administration de la formation : créer un module, déposer ou générer les contenus,
suivre l'avancement de la pipeline et contrôler les artefacts produits. Le deuxième
concerne l'apprenant : accéder aux cours, écouter les audios, consulter les supports et
poser des questions à l'assistant pédagogique. Le troisième concerne la production
automatisée : générer les cours, les audios, les slides et les exercices. Le quatrième
concerne la supervision : vérifier les erreurs, relancer certaines étapes, consulter les
rapports et améliorer les règles.

\begin{figure}[h]
\centering
\fbox{
\begin{minipage}{0.92\linewidth}
\textbf{Diagramme de cas d'utilisation simplifié}

\vspace{0.2cm}
\begin{tabular}{p{3.5cm}p{10cm}}
\textbf{Administrateur} &
Créer une plateforme de formation ; lancer une pipeline RNCP ; suivre les étapes ;
consulter les artefacts ; publier les contenus. \\
\textbf{Apprenant} &
Écouter les cours audio ; consulter les supports ; répondre aux QCM ; poser une
question à l'assistant pédagogique. \\
\textbf{Pipeline IA} &
Générer la base de connaissances ; produire le programme ; rédiger les cours ;
contrôler la qualité ; produire slides et audios. \\
\textbf{Assistant RAG} &
Rechercher les passages pertinents dans l'index documentaire ; générer une réponse
courte et contextualisée. \\
\end{tabular}
\end{minipage}}
\caption{Cas d'utilisation principaux de la plateforme autonome}
\end{figure}

Ce modèle fonctionnel traduit un choix important : la plateforme n'est pas conçue comme
un simple LMS, ni comme un simple chatbot. Elle est conçue comme un système de
production et de diffusion de formation. L'administration, la génération, la diffusion et
l'accompagnement sont pensés ensemble. Cela permet d'éviter une architecture fragmentée
où chaque outil produirait un résultat séparé sans cohérence globale.

\subsection{Principe architectural : un RNCP, un module durable}

Un principe structurant de la solution est le modèle « un RNCP, un module durable ».
Lorsqu'une formation est créée, la pipeline ne produit pas un contenu jetable pour une
promotion unique. Elle produit un module complet, associé à un titre professionnel, qui
peut ensuite être réutilisé pour plusieurs promotions. Ce choix modifie l'économie du
système.

Dans une approche où la pipeline serait relancée pour chaque promotion, le coût de
génération devrait être minimisé à tout prix. Dans le modèle retenu, le coût initial est
amorti sur la durée de vie du module. Il devient donc acceptable d'investir davantage dans
la qualité de la génération, les reviews, les artefacts et les supports, car ces éléments
seront réutilisés. L'enjeu n'est pas de produire le contenu le moins cher possible une seule
fois, mais de produire un module suffisamment fiable pour être exploité dans le temps.

Ce principe a également une conséquence sur l'architecture multi-tenant. Chaque pipeline
crée ou alimente une plateforme dédiée, avec ses contenus, ses fichiers audio et ses
documents. Cette isolation évite de mélanger plusieurs formations, réduit les risques
d'écrasement et facilite la traçabilité. La relation logique devient :

\begin{figure}[h]
\centering
\fbox{
\begin{minipage}{0.88\linewidth}
\centering
\texttt{RNCP} $\rightarrow$ \texttt{Pipeline de génération} $\rightarrow$
\texttt{Plateforme dédiée} $\rightarrow$ \texttt{Module durable} $\rightarrow$
\texttt{Promotions successives}
\end{minipage}}
\caption{Principe d'amortissement : une génération produit un module réutilisable}
\end{figure}

Cette décision répond directement à la problématique économique d'Audit Énergie. La
réduction des coûts ne vient pas seulement du remplacement de certaines tâches humaines
par de l'IA, mais aussi de la réutilisation des productions. Un cours généré, contrôlé,
transformé en audio et enrichi de supports devient un actif pédagogique.

\subsection{Architecture technique générale}

La solution repose sur une architecture web classique complétée par des services IA. Le
backend orchestre les pipelines, stocke l'état des jobs, déclenche les appels aux modèles,
gère les fichiers et expose les routes nécessaires au frontend. Le frontend permet à
l'administrateur de suivre la génération et aux élèves d'accéder aux contenus. Les fichiers
audio et documents sont stockés dans Azure Blob Storage. Le chatbot s'appuie sur Azure
OpenAI et Azure AI Search pour la recherche documentaire.

\begin{figure}[h]
\centering
\fbox{
\begin{minipage}{0.94\linewidth}
\textbf{Architecture logique simplifiée}

\vspace{0.2cm}
\begin{tabular}{p{3.2cm}p{10.2cm}}
\textbf{Frontend React} &
Interface apprenant, dashboard RH, suivi de pipeline, modales d'audit. \\
\textbf{Backend Flask} &
Routes API, orchestration des jobs, contrôle d'accès, stockage des états,
appels aux services IA. \\
\textbf{Base SQLite} &
Configuration des plateformes, jobs de génération, dossiers de cours,
événements et métadonnées. \\
\textbf{Azure Blob Storage} &
Stockage des PDF, audios, documents générés et exports. \\
\textbf{Azure AI Search} &
Indexation documentaire et recherche hybride pour l'assistant pédagogique. \\
\textbf{Modèles IA} &
Génération de contenus, reviews, synthèse vocale, réponses RAG. \\
\end{tabular}
\end{minipage}}
\caption{Architecture technique de la plateforme}
\end{figure}

Ce découpage répond à une contrainte importante : les étapes de génération peuvent être
longues, coûteuses et sujettes à des erreurs réseau ou API. Le backend ne peut donc pas
traiter la pipeline comme une simple requête synchrone. Les jobs doivent être persistés,
leurs statuts doivent être visibles, et certaines étapes doivent pouvoir être relancées sans
reprendre tout le processus.

\subsection{Assistant pédagogique : choix du RAG Azure}

Pour l'assistant pédagogique, le choix retenu est une architecture RAG fondée sur Azure
OpenAI et Azure AI Search. Les documents de formation sont déposés dans un stockage
Azure, indexés par Azure AI Search, puis utilisés comme source documentaire lors des
questions posées par les élèves. Le modèle reçoit la question, l'historique court de la
conversation et les passages récupérés, puis produit une réponse courte.

Ce choix s'inscrit directement dans l'état de l'art. Un prompt seul aurait été insuffisant,
car il ne permet pas de garantir que la réponse vient du cours réellement diffusé. Le
fine-tuning aurait pu adapter le style de l'assistant, mais il n'aurait pas résolu le problème
de mise à jour des documents. Le RAG est donc le choix le plus adapté pour un assistant
dont la connaissance doit suivre les supports de formation.

Le choix d'Azure AI Search est motivé par plusieurs raisons. La première est la recherche
hybride. Les questions des élèves peuvent utiliser un vocabulaire approximatif, tandis que
les documents contiennent des termes métier précis. Une recherche vectorielle simple
risque de négliger certains mots importants ; une recherche par mots-clés seule risque de
manquer les reformulations. La recherche hybride répond à cette double contrainte.

La deuxième raison est l'industrialisation. Azure AI Search fournit un mécanisme d'index,
d'indexer, de stockage documentaire et de séparation par plateforme. Cela correspond
mieux à une plateforme utilisée par de vrais élèves qu'une base vectorielle locale
prototypale. La troisième raison est l'intégration avec Azure OpenAI, qui simplifie
l'utilisation des documents récupérés comme sources de réponse.

\begin{figure}[h]
\centering
\fbox{
\begin{minipage}{0.88\linewidth}
\centering
\texttt{Question élève} $\rightarrow$
\texttt{Azure AI Search} $\rightarrow$
\texttt{Passages pertinents} $\rightarrow$
\texttt{Azure OpenAI} $\rightarrow$
\texttt{Réponse pédagogique courte}
\end{minipage}}
\caption{Chaîne RAG utilisée par l'assistant pédagogique}
\end{figure}

Le prompt système de l'assistant impose plusieurs contraintes : répondre en quelques
phrases, se baser sur le contenu du cours, ne pas inventer d'information présentée comme
issue du cours et adopter un ton pédagogique clair. Cette contrainte de concision est
importante. Dans le contexte d'une formation audio, l'élève pose souvent une question
pour lever un blocage immédiat ; une réponse trop longue peut nuire à l'usage réel.

Le RAG n'est cependant pas utilisé partout. Une décision importante a été de ne pas
l'utiliser pour analyser le REAC lorsque le document entier tient dans la fenêtre de
contexte du modèle. Dans ce cas, le problème n'est pas de retrouver un passage dans un
corpus trop volumineux, mais d'enrichir une source relativement courte pour produire un
contenu pédagogique dense. Cette distinction entre problème de récupération et problème
de densité est un apport méthodologique important du projet.

\subsection{Pipeline de formation : de la source au module}

La deuxième brique centrale est la pipeline de génération. Son rôle est de transformer un
référentiel ou un ensemble de sources en une formation structurée. Les premières versions
reposaient sur une logique plus directe : extraction du contenu, découpage en journées,
génération de textes longs, puis transformation audio. Cette approche a permis de valider
la faisabilité, mais elle a montré ses limites pour la traçabilité, la cohérence et la qualité.

La pipeline actuelle est organisée comme une succession d'étapes spécialisées :
initialisation du module, récupération ou dépôt des sources, enrichissement de la base de
connaissances, génération du programme, découpage journalier, génération structurée des
cours, reviews, création des supports, préparation audio et publication. Chaque étape
produit un état ou un artefact exploitable.

\begin{figure}[h]
\centering
\fbox{
\begin{minipage}{0.94\linewidth}
\centering
\texttt{RNCP / REAC}
$\rightarrow$ \texttt{Knowledge base enrichie}
$\rightarrow$ \texttt{Programme global}
$\rightarrow$ \texttt{Journées}
$\rightarrow$ \texttt{Plan JSON}
$\rightarrow$ \texttt{Sections}
$\rightarrow$ \texttt{Reviews}
$\rightarrow$ \texttt{Slides + audio}
$\rightarrow$ \texttt{Module publié}
\end{minipage}}
\caption{Pipeline générale de production d'un module de formation}
\end{figure}

Cette architecture se distingue d'une génération directe par prompt unique. Le modèle ne
porte pas toute la responsabilité en une seule sortie. Il intervient dans des tâches plus
restreintes : extraire, enrichir, planifier, rédiger, corriger, vérifier, reformuler ou produire
un support. Cette décomposition réduit les dérives et rend les erreurs plus localisables.

\subsection{Enrichissement de la source pédagogique}

Le REAC constitue une source officielle, mais il n'est pas toujours suffisant pour produire
des heures de cours exploitables. Il décrit les compétences, les attendus, les savoir-faire
et les conditions d'exercice, mais il ne fournit pas nécessairement assez d'exemples, de
cas pratiques, de vocabulaire métier ou de situations pédagogiques pour alimenter une
formation audio longue.

La pipeline introduit donc une étape d'enrichissement. À partir des compétences extraites
du REAC, le système génère une base de connaissances pédagogique : définitions, études
de cas, pièges fréquents, vocabulaire, contexte terrain et liens entre compétences. Cette
base n'est pas destinée directement à l'élève ; elle sert de matière première pour la
génération des cours.

Cette étape répond à une limite importante des approches naïves. Demander au modèle
de produire plusieurs centaines de milliers de mots de formation à partir d'un référentiel
relativement court crée un risque de dilution. Le modèle risque de répéter les mêmes
idées ou de combler les vides avec des généralités. L'enrichissement vise donc à augmenter
la densité pédagogique avant la génération finale.

Le choix retenu n'est pas un RAG sur le REAC, car le REAC tient dans la fenêtre de
contexte du modèle. Le problème n'est pas la recherche d'information, mais la production
de matière pédagogique structurée. Le RAG reste pertinent pour l'assistant élève ou pour
de futurs corpus externes volumineux ; il ne constitue pas la bonne réponse à ce stade de
la pipeline.

\subsection{Plan JSON comme contrat pédagogique}

Le cœur de l'approche originale est le passage d'une génération directe à une génération
fondée sur un plan JSON. Ce plan joue le rôle de contrat pédagogique. Il décrit les cours
de la journée, les sections, les objectifs, les budgets de mots, les transitions, les moments
visualisables et les ancrages de slides. Le texte final n'est plus le premier objet produit :
il est la conséquence d'un plan explicite.

Ce choix répond à plusieurs problèmes identifiés dans les versions précédentes. D'abord,
il améliore la cohérence. Une journée de formation longue doit suivre une progression.
Sans plan explicite, les introductions peuvent se répéter, les transitions devenir
artificielles et les conclusions fermer le mauvais sujet. Ensuite, il améliore la traçabilité :
si une section dévie, on peut comparer le texte au plan prévu. Enfin, il rend possible la
génération de supports alignés, car les moments pédagogiques importants sont identifiés
avant la création des slides.

\begin{figure}[h]
\centering
\fbox{
\begin{minipage}{0.88\linewidth}
\textbf{Structure logique du plan JSON}

\vspace{0.2cm}
\texttt{Journée}
$\rightarrow$ \texttt{Cours 1..7}
$\rightarrow$ \texttt{Opening / Parts / Conclusion}
$\rightarrow$ \texttt{Teaching beats}
$\rightarrow$ \texttt{Slide anchors}
\end{minipage}}
\caption{Le plan JSON comme contrat pédagogique et visuel}
\end{figure}

Le plan JSON permet également de séparer le langage interne de la parole du formateur.
Le code peut manipuler des notions de cours, sections, budgets, blocs audio ou anchors.
Mais l'élève ne doit pas entendre ces contraintes techniques. Le texte final doit rester
naturel : parler de thème, de point, de partie, de reprise après pause ou de fil conducteur.
Cette séparation est essentielle pour éviter que la mécanique de la pipeline apparaisse
dans le discours pédagogique.

\subsection{Génération par sections et contrôle des budgets}

Une fois le plan établi, la génération se fait section par section. Chaque section dispose
d'un rôle, d'un objectif, d'un budget et d'un périmètre. Cette décision réduit la charge
cognitive imposée au modèle. Au lieu de générer plusieurs heures de cours en conservant
toutes les contraintes en mémoire, le modèle rédige une unité limitée.

Le budget de mots est une contrainte centrale, car le texte doit ensuite être transformé
en audio. Une génération trop courte donne une formation pauvre ou déséquilibrée. Une
génération trop longue dépasse les créneaux prévus, complique le découpage et peut
produire des audios inutilisables. La pipeline donne donc au serveur un rôle d'autorité sur
les budgets. Les budgets ne sont pas seulement des consignes textuelles données au
modèle ; ils sont contrôlés et recalibrés par le système.

Cette approche se distingue d'un simple prompt qui demanderait au modèle de respecter
une durée. Les modèles de langage ne maîtrisent pas précisément la correspondance entre
mots, durée audio et rythme de synthèse vocale. Il faut donc traiter le budget comme un
invariant logiciel, avec des étapes de vérification et de correction. Cette décision rend la
pipeline plus robuste, même si elle introduit un coût supplémentaire en appels modèle et
en logique d'orchestration.

\subsection{Reviews ciblées et auditabilité}

La qualité de la génération est contrôlée par plusieurs reviews spécialisées. L'objectif
n'est pas de demander à un modèle de « rendre le texte meilleur » de manière générale,
mais d'isoler les risques. Une review vérifie l'adhérence au plan ; une autre ajuste le
volume ; une micro-review éthique traite localement les formulations sensibles ; une
review d'humanisation améliore l'oralité ; une conformité finale vérifie les règles les plus
critiques.

Cette séparation est un apport important de l'approche. Dans les versions plus anciennes,
les reviews globales mélangeaient plusieurs dimensions : structure, style, conformité,
oralité et parfois budget. Une correction pouvait alors résoudre un problème tout en en
créant un autre. En spécialisant les reviews, chaque étape dispose d'une responsabilité
plus claire.

L'auditabilité repose sur la persistance d'artefacts intermédiaires. La pipeline conserve
notamment le plan, les sections brutes, les scripts assemblés, les rapports de review, les
versions corrigées, le plan audio et les artefacts de slides. Ces éléments transforment la
génération IA en processus vérifiable. Si une erreur apparaît dans le cours final, on peut
chercher à quelle étape elle a été introduite.

\begin{figure}[h]
\centering
\fbox{
\begin{minipage}{0.92\linewidth}
\textbf{Séquence simplifiée de génération d'un cours}

\vspace{0.2cm}
\texttt{Plan JSON}
$\rightarrow$ \texttt{Sections brutes}
$\rightarrow$ \texttt{Audit adhérence}
$\rightarrow$ \texttt{Budget}
$\rightarrow$ \texttt{Micro-conformité}
$\rightarrow$ \texttt{Humanisation}
$\rightarrow$ \texttt{Conformité finale}
$\rightarrow$ \texttt{Audio-plan}
\end{minipage}}
\caption{Séquence de contrôle appliquée aux contenus générés}
\end{figure}

L'interface d'administration rend ces artefacts consultables. Les étapes de la roadmap
peuvent être ouvertes dans des modales d'audit. Lorsqu'une micro-conformité éthique
modifie un passage, l'interface peut afficher le texte original, le texte corrigé, la règle
concernée et la justification. Cette fonctionnalité est importante : l'auditabilité ne doit
pas exister seulement dans les logs techniques, elle doit être accessible à l'utilisateur qui
pilote la formation.

\subsection{Slides anchor-first et cohérence pédagogique}

La génération des supports visuels pose un problème spécifique. Une solution simple
consisterait à générer le texte du cours, puis à demander au modèle d'en extraire les
points principaux pour créer des slides. Cette approche est facile à mettre en place, mais
elle produit souvent des supports génériques. Les éléments les plus fréquents dans un
texte ne sont pas toujours les plus importants à visualiser.

L'approche retenue est une génération anchor-first. Les moments pédagogiques à
visualiser sont décidés dès le plan JSON sous forme de teaching beats et de slide anchors.
Une slide peut correspondre à une définition, une checklist, un piège, un exemple, une
comparaison, une méthode ou un cas pratique. Le texte couvre ensuite ces moments
naturellement, sans verbaliser la présence des slides ou des templates.

Cette stratégie présente deux avantages. D'abord, elle aligne les supports sur l'intention
pédagogique plutôt que sur une analyse opportuniste du texte final. Ensuite, elle rend les
slides traçables : chaque support peut être rattaché à un anchor prévu dans le plan et à
un passage du script. La slide n'est donc pas une décoration ajoutée après coup, mais une
partie du dispositif pédagogique.

\subsection{Audio et transformation TTS}

La sortie principale de la formation reste l'audio. Les textes générés doivent donc être
compatibles avec une synthèse vocale. Cette contrainte influence toute la pipeline. Le
style doit être oral, les phrases doivent être prononçables, les transitions doivent être
naturelles et la durée doit correspondre aux blocs prévus.

Les premières expérimentations avec des outils de synthèse vocale, comme ElevenLabs,
ont montré que la qualité de la voix ne suffit pas. Un bon moteur TTS peut rendre un
texte agréable à écouter, mais il ne corrige pas un mauvais script. La pipeline traite donc
la qualité audio en amont, par la rédaction TTS-ready, le contrôle du budget, l'humanisation
orale et le plan audio.

La génération audio reste toutefois une zone à évaluer finement. La correspondance entre
nombre de mots et durée réelle dépend de la voix, du rythme, des pauses et du moteur
utilisé. Les ratios de mots par minute employés dans la pipeline doivent donc être
considérés comme des paramètres de calibration, non comme des résultats validés. Une
évaluation plus rigoureuse doit mesurer l'écart entre durée prévue et durée réelle sur un
échantillon d'audios générés.

\subsection{Résultats obtenus}

Les résultats du projet peuvent être analysés selon plusieurs dimensions. Sur le plan
fonctionnel, la plateforme permet de diffuser des contenus de formation à de vrais élèves,
de gérer des modules de cours, de déposer des documents et d'interagir avec un assistant
pédagogique. Le projet dépasse donc le stade d'une expérimentation isolée : il est intégré
à un environnement applicatif utilisé.

Sur le plan de la génération, la pipeline a évolué d'une génération directe ou slotée vers
une architecture structurée. La présence d'un plan JSON, de sections, de reviews et
d'artefacts intermédiaires améliore la capacité à contrôler le résultat. L'amélioration la
plus importante n'est pas seulement qualitative ; elle est méthodologique. Le système ne
produit plus uniquement un texte final, il produit une chaîne de preuves.

Sur le plan de l'accompagnement, l'assistant RAG permet de répondre aux questions à
partir des documents indexés, ce qui réduit les limites d'un cours audio préenregistré. Il
ne remplace pas toutes les formes d'interaction humaine, mais il fournit un premier niveau
de réponse contextualisée, disponible pendant la formation.

Sur le plan organisationnel, le modèle « un RNCP, un module durable » transforme la
production en investissement réutilisable. La pipeline peut coûter plus cher qu'une
génération ponctuelle, mais ce coût est amorti sur les promotions successives. Cette
logique est cohérente avec les contraintes économiques d'une petite structure.

\begin{table}[h]
\centering
\small
\begin{tabular}{|p{3.3cm}|p{4.4cm}|p{4.4cm}|}
\hline
\textbf{Dimension} & \textbf{Avant la refonte structurée} & \textbf{Après l'approche proposée} \\
\hline
Structure du cours &
Texte long généré par créneaux ou blocs, avec contrôle limité de la progression. &
Plan JSON explicite, sections, objectifs, budgets et anchors pédagogiques. \\
\hline
Qualité &
Reviews plus globales, erreurs détectées tardivement. &
Reviews ciblées : adhérence au plan, budget, micro-conformité, humanisation, conformité finale. \\
\hline
Traçabilité &
Sortie principalement textuelle, difficile à auditer. &
Artefacts intermédiaires, rapports, diff avant/après et modales d'audit. \\
\hline
Slides &
Risque de génération après coup par analyse du texte final. &
Approche anchor-first pilotée par le plan pédagogique. \\
\hline
Assistant élève &
Cours audio peu interactif. &
Assistant RAG contextualisé sur les documents de formation. \\
\hline
Modèle économique &
Production fortement dépendante d'interventions humaines répétées. &
Module durable réutilisable, coût de génération amorti sur plusieurs promotions. \\
\hline
\end{tabular}
\caption{Comparaison synthétique avant/après l'approche proposée}
\end{table}

\subsection{Comparaison avec les approches de l'état de l'art}

Par rapport au prompt engineering simple, l'approche proposée ajoute une architecture
autour du modèle. Les prompts restent nécessaires, mais ils ne portent plus seuls la
qualité du système. Ils sont spécialisés par rôle et intégrés dans une pipeline qui conserve
des artefacts et impose des invariants.

Par rapport au fine-tuning, la solution privilégie la flexibilité documentaire. Les contenus
de formation peuvent changer, être indexés, enrichis ou régénérés sans réentraîner un
modèle. Le fine-tuning pourrait éventuellement être utile plus tard pour stabiliser le style
de l'assistant ou certaines reviews, mais il n'est pas le mécanisme principal.

Par rapport à un RAG simple, le choix d'Azure AI Search permet une approche plus
industrialisée : indexation, recherche hybride, séparation des corpus par plateforme et
intégration avec Azure OpenAI. Cette solution reste dépendante de la qualité des documents
indexés, mais elle correspond mieux à un usage multi-modules qu'un prototype local.

Par rapport à une génération directe de cours, la pipeline structurée apporte une meilleure
maîtrise des contenus longs. Le plan JSON, les sections et les reviews permettent de
localiser les erreurs et de discuter la qualité de manière plus objective. Cette approche
est plus complexe et plus coûteuse, mais elle est cohérente avec le modèle de module
durable.

Par rapport à une solution human-in-the-loop classique, l'approche réduit l'intervention
humaine récurrente sans supprimer toute supervision. L'humain intervient davantage pour
piloter, valider, auditer et améliorer le système. Cette position est importante : la solution
vise l'autonomie opérationnelle, pas l'absence de responsabilité.

\subsection{Évaluations à consolider}

Pour rendre l'étude pleinement scientifique, plusieurs évaluations doivent être réalisées
ou complétées. Certaines peuvent être menées à partir des logs et artefacts existants ;
d'autres nécessitent de construire un protocole dédié.

La première évaluation concerne le RAG. Il faut constituer un jeu de questions
représentatives des élèves, identifier pour chaque question les passages sources attendus,
puis mesurer la récupération et la réponse. Les métriques pertinentes sont la précision@k,
le rappel@k, le rang du premier passage pertinent, le taux de réponses supportées par les
sources et le taux d'hallucinations. Une comparaison entre prompt seul et RAG permettrait
de démontrer l'apport réel de l'ancrage documentaire.

La deuxième évaluation concerne les variantes de recherche. Lorsque la configuration le
permet, il faut comparer recherche vectorielle simple, recherche hybride et recherche
hybride avec reclassement sémantique. Cette comparaison est importante pour justifier
le choix d'Azure AI Search autrement que par des arguments d'intégration.

La troisième évaluation concerne la pipeline de génération. Il faut comparer un contenu
produit par génération directe et un contenu produit par la pipeline structurée. Les critères
peuvent être le respect du plan, le nombre de répétitions, le respect du budget de mots,
la qualité orale, le nombre de corrections nécessaires et la cohérence entre cours, slides
et QCM.

La quatrième évaluation concerne les reviews. Il faut mesurer combien de problèmes sont
détectés par chaque review, combien de corrections sont appliquées, combien sont
rejetées et combien produisent un effet indésirable. Cela permettrait de savoir quelles
reviews apportent réellement de la valeur et lesquelles doivent être améliorées.

La cinquième évaluation concerne l'impact économique et organisationnel. Il faut mesurer
le temps nécessaire pour produire une journée ou un module avant et après automatisation,
le nombre d'interventions humaines restantes, le coût API et TTS, ainsi que la capacité à
réutiliser le module sur plusieurs promotions. Ces mesures sont nécessaires pour relier la
solution à la problématique initiale.

\subsection{Conditions d'application}

L'approche proposée est adaptée à des formations structurées, répétables et fondées sur
un référentiel relativement stable. Elle est particulièrement pertinente lorsque le contenu
peut être préparé en amont, transformé en audio et réutilisé pour plusieurs promotions.
Elle convient donc bien aux titres professionnels ou aux formations standardisées.

Elle est moins adaptée à des formations très interactives, fortement dépendantes du débat,
de la pratique immédiate ou de l'improvisation pédagogique. Dans ces cas, un formateur
humain reste difficile à remplacer. La plateforme peut alors servir de support ou de
complément, mais pas nécessairement de substitut principal.

L'approche suppose également un corpus de départ de qualité. Si le référentiel est trop
pauvre, incomplet ou ambigu, la génération risque de produire des contenus génériques.
L'enrichissement de la base de connaissances réduit ce risque, mais ne peut pas inventer
une expertise métier absente des sources ou de la supervision.

Enfin, l'approche nécessite une infrastructure capable de gérer des traitements longs :
stockage durable, jobs persistés, reprise après erreur, logs, suivi des coûts et supervision
des services externes. Une pipeline IA longue ne peut pas être traitée comme une simple
fonction appelée à la demande.

\subsection{Limites de la contribution}

La première limite concerne l'évaluation quantitative. La plateforme fonctionne avec de
vrais élèves et la pipeline produit des artefacts exploitables, mais certaines métriques
doivent encore être consolidées : temps gagné, coût complet, satisfaction des élèves,
fidélité du RAG, précision de la durée audio et impact sur les résultats pédagogiques.
Ces métriques sont identifiées, mais elles doivent faire l'objet d'un protocole plus
systématique.

La deuxième limite concerne la dépendance aux services externes. La solution s'appuie
sur des modèles de langage, des services Azure, des outils de synthèse vocale et des API.
Ces services peuvent évoluer, changer de coût ou rencontrer des limites de débit. La
robustesse du système dépend donc aussi de la gestion des erreurs, des retries, de la
surveillance et de la possibilité de remplacer certaines briques.

La troisième limite concerne la qualité pédagogique profonde. Les reviews automatiques
peuvent détecter certaines erreurs, mais elles ne remplacent pas entièrement le jugement
d'un expert pédagogique. Elles améliorent l'observabilité et réduisent certains risques,
mais elles doivent être évaluées et supervisées.

La quatrième limite concerne la rigidité du modèle. La pipeline est adaptée à un format
de journée structuré, avec des blocs audio et une logique de module durable. Cette
structure est cohérente avec le besoin actuel, mais elle peut être moins adaptée à des
formations plus courtes, plus libres ou plus interactives. Une évolution future devra
permettre de paramétrer davantage les formats.

\subsection{Pistes d'amélioration}

Plusieurs pistes d'amélioration se dégagent. La première est la mise en place d'un banc
d'évaluation RAG. Il faudrait créer un jeu de questions annotées, comparer plusieurs
configurations de recherche et conserver des métriques dans le temps. Cela permettrait
de détecter les régressions lorsque les documents, les prompts ou les modèles changent.

La deuxième piste est l'évaluation automatique et humaine des contenus générés. Une
grille de relecture pourrait mesurer la fidélité au programme, la clarté, l'oralité, la
cohérence des slides et la qualité des QCM. Un échantillon de cours pourrait être évalué
par un humain, puis comparé aux reviews automatiques.

La troisième piste est l'amélioration de la calibration audio. Il faudrait mesurer
systématiquement les durées réelles des fichiers générés, comparer ces durées aux
budgets prévus et ajuster les ratios par voix ou par moteur TTS. Cela permettrait de
remplacer une calibration empirique par une calibration mesurée.

La quatrième piste est l'automatisation complète du provisionnement multi-tenant.
Aujourd'hui, certains éléments d'infrastructure peuvent encore nécessiter une intervention
manuelle. À terme, la création d'un module devrait provisionner automatiquement les
ressources nécessaires : containers, index, configuration de recherche et droits associés.

La cinquième piste est l'amélioration du rôle humain. Plutôt que de relire tout le contenu,
l'interface pourrait proposer une validation ciblée : sections signalées comme risquées,
corrections importantes, slides sans source forte, réponses RAG à faible confiance ou
écarts de durée audio. L'humain interviendrait alors sur les points réellement incertains.

\subsection{Apports professionnels et montée en compétences}

Ce projet a permis de développer des compétences qui dépassent la simple utilisation
d'outils d'intelligence artificielle. La première compétence est l'analyse d'un besoin métier.
Le problème initial était économique et organisationnel : rendre viable une activité de
formation pour une petite structure. La solution technique n'a de sens que parce qu'elle
répond à cette contrainte.

La deuxième compétence est l'architecture logicielle. Il a fallu concevoir une plateforme
capable de gérer des traitements longs, des fichiers, des jobs, des statuts, des services
externes et plusieurs modules de formation. La persistance, la reprise après erreur et la
séparation des responsabilités sont devenues des sujets centraux.

La troisième compétence est l'ingénierie des systèmes IA. Le projet a demandé de choisir
quand utiliser un prompt, quand utiliser un RAG, quand enrichir une source, quand
structurer une génération et quand ajouter une review. L'enjeu était de ne pas appliquer
l'IA comme une solution générique, mais de choisir la bonne technique pour le bon
problème.

La quatrième compétence est la prise de recul. Les premières solutions ont fonctionné
partiellement, mais elles ont révélé de nouvelles limites. Le passage à une pipeline
auditable est le résultat de ces itérations. Cette évolution montre l'importance de mesurer,
de comparer et de remettre en question une architecture lorsqu'elle devient difficile à
contrôler.

Enfin, le projet a renforcé une compétence de pilotage produit. Une plateforme autonome
ne se juge pas seulement par son code. Elle doit être utilisable, compréhensible, maintenable
et économiquement pertinente. L'étude originale montre ainsi comment un besoin
d'entreprise a conduit à une solution technique structurée, puis à une réflexion plus large
sur l'évaluation et les limites de l'automatisation par IA.

% ====================== CONCLUSION ======================
\clearpage
\phantomsection
\addcontentsline{toc}{section}{Conclusion générale}
\section*{Conclusion générale}

Ce mémoire est parti d'un besoin concret et profondément économique : permettre à une
petite structure, Audit Énergie, de développer une activité de formation externe sans dépendre
du coût récurrent d'un formateur en direct. Cette contrainte initiale a guidé l'ensemble du
travail. À chaque étape, une solution a été essayée, ses limites ont été observées, puis une
réponse plus complète a été conçue : du cours en direct au cours audio préparé à l'avance, puis
aux essais de synthèse vocale, et enfin à une plateforme autonome combinant une pipeline de
génération de contenus et un assistant pédagogique contextualisé.

L'apport central du travail n'est pas l'utilisation isolée d'un modèle de langage ou d'un service
de synthèse vocale, mais leur \emph{orchestration} dans une chaîne fiable et observable. La
pipeline structurée, fondée sur un plan explicite, une génération section par section, des
reviews ciblées et des artefacts intermédiaires, transforme une sortie d'IA opaque en un
processus auditable. L'assistant RAG, ancré sur les documents réels de la formation, prolonge
l'expérience du cours en répondant aux questions des élèves sans tomber dans les dérives d'un
chatbot généraliste. Ensemble, ces deux briques répondent à la problématique posée :
automatiser fortement le modèle de formation tout en préservant la qualité pédagogique, la
traçabilité des contenus et la pertinence des réponses.

Ce travail comporte des limites assumées. Plusieurs indicateurs — temps réellement gagné,
coût complet, fidélité du RAG, précision de la calibration audio, satisfaction des élèves —
restent à consolider par un protocole d'évaluation plus systématique. La solution dépend par
ailleurs de services externes dont le coût et la disponibilité peuvent évoluer, et la supervision
humaine demeure nécessaire sur les points les plus sensibles. Ces limites ne remettent pas en
cause l'approche : elles tracent les pistes d'évolution, notamment la mise en place d'un banc
d'évaluation RAG, l'évaluation croisée automatique et humaine des contenus générés, et une
calibration audio mesurée plutôt qu'empirique.

Sur le plan professionnel, ce projet a permis une véritable montée en compétences :
l'analyse d'un besoin métier, la conception d'une architecture logicielle capable de gérer des
traitements longs et reprenables, l'ingénierie des systèmes d'IA (savoir quand recourir à un
prompt, à un RAG, à un enrichissement de source ou à une génération structurée) et, surtout,
la prise de recul nécessaire pour remettre en question une solution lorsqu'elle devient
difficile à contrôler. Au-delà de la plateforme elle-même, ce mémoire illustre une conviction
construite tout au long du travail : la valeur de l'intelligence artificielle générative en
entreprise ne tient pas à la puissance brute des modèles, mais à la rigueur de l'architecture
qui les encadre.

% ====================== BIBLIOGRAPHIE ===============================
\clearpage
\phantomsection
\addcontentsline{toc}{section}{Bibliographie}
\begin{thebibliography}{99}

\bibitem{brown2020language} T. B. Brown, B. Mann, N. Ryder \emph{et al.}, \emph{Language Models are Few-Shot Learners}, Advances in Neural Information Processing Systems (NeurIPS), 2020.

\bibitem{hu2021lora} E. J. Hu, Y. Shen, P. Wallis \emph{et al.}, \emph{LoRA: Low-Rank Adaptation of Large Language Models}, arXiv:2106.09685, 2021.

\bibitem{lewis2020retrieval} P. Lewis, E. Perez, A. Piktus \emph{et al.}, \emph{Retrieval-Augmented Generation for Knowledge-Intensive NLP Tasks}, Advances in Neural Information Processing Systems (NeurIPS), 2020.

\bibitem{ji2023hallucination} Z. Ji, N. Lee, R. Frieske \emph{et al.}, \emph{Survey of Hallucination in Natural Language Generation}, ACM Computing Surveys, vol.~55, n\textsuperscript{o}~12, 2023.

\bibitem{gao2023retrieval} Y. Gao, Y. Xiong, X. Gao \emph{et al.}, \emph{Retrieval-Augmented Generation for Large Language Models: A Survey}, arXiv:2312.10997, 2023.

\bibitem{microsoft2026foundryrag} Microsoft, \emph{Retrieval Augmented Generation (RAG) in Azure AI Search}, Microsoft Learn, documentation Azure AI Search, 2026.

\bibitem{microsoft2026hybridsearch} Microsoft, \emph{Hybrid search in Azure AI Search}, Microsoft Learn, documentation Azure AI Search, 2026.

\bibitem{microsoft2026semanticranker} Microsoft, \emph{Semantic ranking in Azure AI Search}, Microsoft Learn, documentation Azure AI Search, 2026.

\bibitem{es2023ragas} S. Es, J. James, L. Espinosa-Anke, S. Schockaert, \emph{RAGAS: Automated Evaluation of Retrieval Augmented Generation}, arXiv:2309.15217, 2023.

\bibitem{saadfalcon2023ares} J. Saad-Falcon, O. Khattab, C. Potts, M. Zaharia, \emph{ARES: An Automated Evaluation Framework for Retrieval-Augmented Generation Systems}, arXiv:2311.09476, 2023.

\bibitem{liu2023evaluating} N. F. Liu, T. Zhang, P. Liang, \emph{Evaluating Verifiability in Generative Search Engines}, Findings of EMNLP, 2023.

\bibitem{wei2022chain} J. Wei, X. Wang, D. Schuurmans \emph{et al.}, \emph{Chain-of-Thought Prompting Elicits Reasoning in Large Language Models}, Advances in Neural Information Processing Systems (NeurIPS), 2022.

\bibitem{zhou2022least} D. Zhou, N. Schärli, L. Hou \emph{et al.}, \emph{Least-to-Most Prompting Enables Complex Reasoning in Large Language Models}, arXiv:2205.10625, 2022.

\bibitem{yao2023tree} S. Yao, D. Yu, J. Zhao \emph{et al.}, \emph{Tree of Thoughts: Deliberate Problem Solving with Large Language Models}, Advances in Neural Information Processing Systems (NeurIPS), 2023.

\bibitem{shinn2023reflexion} N. Shinn, F. Cassano, A. Gopinath, K. Narasimhan, S. Yao, \emph{Reflexion: Language Agents with Verbal Reinforcement Learning}, Advances in Neural Information Processing Systems (NeurIPS), 2023.

\bibitem{zheng2023judging} L. Zheng, W.-L. Chiang, Y. Sheng \emph{et al.}, \emph{Judging LLM-as-a-Judge with MT-Bench and Chatbot Arena}, Advances in Neural Information Processing Systems (NeurIPS), 2023.

\end{thebibliography}

\end{document}
